\documentclass[11pt,a4paper]{article}

\usepackage[margin=1in]{geometry}
\usepackage{amsmath,amssymb,amsthm}
\usepackage{booktabs}
\usepackage[expansion=false]{microtype}
\usepackage{enumitem}
\usepackage{xcolor}
\usepackage[hidelinks]{hyperref}
\usepackage{url}

\theoremstyle{definition}
\newtheorem{definition}{Definition}
\newtheorem{observation}{Observation}
\newtheorem{attack}{Attack}
\theoremstyle{plain}
\newtheorem{proposition}{Proposition}
\newtheorem{claim}{Claim}

\newcommand{\Gb}{\mathsf{Gb}}
\newcommand{\En}{\mathsf{En}}
\newcommand{\Ev}{\mathsf{Ev}}
\newcommand{\De}{\mathsf{De}}
\newcommand{\lsbf}{\mathsf{lsb}}
\newcommand{\bit}{\{0,1\}}
\newcommand{\Hh}{H}

\title{\bf Free-XOR Garbling Is Not Binding Against the Garbler\\[0.4em]
\large What the On-Chain Label Reveal Buys in BitVM-Style Protocols}

\author{Aria Naraghi\thanks{Independent. Correspondence: \texttt{arianari97@gmail.com}.
Artifacts: \url{https://github.com/aryaethn/garbler-binding}}}

\date{August 2026 \\[0.3em] \small Preprint. Comments welcome.}

\begin{document}
\maketitle

\begin{abstract}
BitVM-style protocols move SNARK verification off the Bitcoin chain by having a
prover garble a verifier circuit and having challengers evaluate it. Soundness
rests on the prover being unable to make an honest evaluation accept without a
valid witness. We observe that no standard garbling security property gives
this: correctness, privacy and authenticity are all stated against adversaries
who do \emph{not} hold the encoding information, and in this setting the
adversary is the garbler, who holds all of it. Every deployed and proposed
design closes the gap the same way, by revealing the garbled input labels
on-chain, which is why on-chain cost is linear in the size of the proof.

We isolate the missing property, which we call \emph{evaluation-binding}, and
show it is not free. Under the native decoding of the half-gates scheme it fails
in two trials. Under BitVM3-core's decoding commitment it is unobservable,
because only the false output label is committed. Stated in the natural way it
is false for a trivial reason. Repaired, its security is at most $2^{k/2}$
rather than $2^{k}$ whenever the circuit admits an affine split reachable from
the output, which free-XOR makes structurally likely: we give a working birthday
attack, measured at scaling exponent $1.02$ over $k \in [24,44]$, that produces
label vectors that encode no input yet decode to \textsc{True}.

We then ask whether the circuit actually deployed has the structure. Using a
custom evaluation mode over the BN254 Groth16 verifier of the BitVM Alliance
implementation ($1.03 \times 10^{10}$ gates, $2.71 \times 10^{9}$ non-free,
$3302$ input wires), we show the output's linear frontier is a single atom and
that none of the $638{,}658$ non-free gates in the circuit's tail has separably
controllable operands. The deployed system is therefore \emph{not} vulnerable,
but it is safe by accident: the property that saves it is a consequence of
Groth16's final equality check, is stated nowhere, and is not obviously
preserved under a change of proof system. We contribute the definition, the
attacks, an exploit, and a tool that decides the structural condition.
\end{abstract}

\section{Introduction}

Bitcoin cannot verify a SNARK. A Groth16 verifier~\cite{Groth16} compiles to
roughly a gigabyte of Bitcoin Script and a block holds four megabytes. The
BitVM line of work~\cite{BitVM1,BitVM2,BitVM3} answers this optimistically: do
not verify on-chain, but arrange that an incorrect claim can be
\emph{disproved} on-chain cheaply.

BitVM3~\cite{BitVM3} realises the fraud proof with garbled circuits. An operator
garbles the SNARK verifier once, publishes the garbled circuit off-chain, and
posts a transaction whose witness reveals the garbled encoding of its claimed
proof. Any party may evaluate the circuit on that encoding; if the proof is
invalid the evaluation yields a designated label $L^{*}$, and revealing $L^{*}$
on-chain spends the operator's deposit. The reported cost is roughly \$9 total
with a fraud proof under \$0.20, a thousandfold improvement on BitVM2.

The on-chain reveal is expensive and it is the reason on-chain cost is $O(|\pi|)$
in the proof size. Three separate lines of recent work exist to manage its
consequences. Eagen's Glock~\cite{Glock} identifies why it works at all,
observing that projective signatures compose with garbled input encodings in
what he calls an ``auspicious harmony,'' and then makes the identification
literal by \emph{defining} input labels to be signature shares. Duty-Free
Bits~\cite{DutyFreeBits} exists because that identification forces projectivity:
its authors state plainly that projectivity is mandated by the on-chain
authentication mechanism. Mosaic~\cite{Mosaic} reduces the on-chain footprint so
it no longer scales with the number of cut-and-choose copies, and opens by
noting that the verifier needs the prover's input labels, which the prover must
post on chain.

\subsection*{The gap}

What does the on-chain reveal actually buy, cryptographically?

BitVM3's soundness argument answers this precisely, in its own words: a
malicious prover posting the assertion reveals, \emph{by extractability} of the
signature scheme, the encoding $\En(e,\pi)$; and then \emph{by correctness} of
the garbling scheme, evaluating on that encoding reproduces the verifier's
verdict. The second step depends on the first. Correctness of a garbling scheme
is a statement about valid encodings and says nothing about $\Ev(F, L)$ for $L$
outside the encoding space. Extractability is what supplies the hypothesis that
correctness requires.

So the on-chain reveal purchases exactly one thing: the guarantee that the
labels the evaluator acts on are well-formed. Remove it and the soundness
argument does not weaken, it loses its first line.

One might hope authenticity covers the residue. It does not, and the definition
says so. In the formulation of Zahur, Rosulek and Evans~\cite{ZRE15}, the
authenticity adversary is given $(F, X)$ \emph{alone}: not the encoding
information $e$, not the decoding information $d$, not the global offset $R$.
The garbler holds all three. Authenticity is silent about the garbler by
construction, and correctly so, since it was formulated for a different threat
model~\cite{BHR12}.

\subsection*{Contributions}

\begin{enumerate}[leftmargin=*,itemsep=2pt]
  \item \textbf{A definition.} We isolate \emph{evaluation-binding}
    (Definition~\ref{def:evb}): no PPT adversary holding $(F, e, d)$ can produce
    an input label vector on which honest evaluation decodes to \textsc{True}
    unless the circuit is satisfiable. This is the property BitVM-style
    protocols need and obtain, today, only through the on-chain reveal. It is
    not implied by correctness, privacy or authenticity.

  \item \textbf{Five attacks} (Section~\ref{sec:attacks}). Native half-gates
    decoding admits forgery in two expected trials. BitVM3-core commits to the
    false output label only, making success unobservable off-chain. The
    encoding-form statement of the property is false for a trivial free-XOR
    reason. Even fully repaired, the property has a $2^{k/2}$ ceiling. And the
    existing soundness proof admits no partial repair.

  \item \textbf{A working exploit} (Section~\ref{sec:impl}). A faithful,
    $k$-parameterised implementation of the half-gates scheme, validated
    exhaustively against plaintext evaluation, in which the birthday attack
    produces label vectors that are the encoding of \emph{no} input yet decode
    to \textsc{True} under committed decoding. Measured scaling exponent $1.02$
    against a birthday prediction of $1.00$, over $k \in [24,44]$; speedup over
    naive search of $3259\times$ already at $k=28$.

  \item \textbf{A measurement of the deployed circuit}
    (Section~\ref{sec:measure}). A custom evaluation mode over the BitVM
    Alliance BN254 Groth16 verifier~\cite{GSV} establishes that its output's
    linear frontier is a single atom and that none of the $638{,}658$ non-free
    gates in its tail has separably controllable operands. The attack does not
    apply to it.

  \item \textbf{A design requirement.} The deployed circuit is safe because
    Groth16 ends in an equality check that aggregates the whole pairing result.
    Nobody chose that. We argue evaluation-binding should become an explicit,
    checked requirement, and we release the tool that checks it.
\end{enumerate}

We emphasise what this paper does \emph{not} claim. We do not break any deployed
system. The deployed system is protected by the on-chain reveal, which makes
evaluation-binding moot; and even without it, we measured the circuit and it
resists our attack. The claim is that the guarantee is unstated, the margin is
smaller than the parameter suggests, and the protection is incidental.

\section{Preliminaries}

\subsection{Garbling schemes}

A garbling scheme~\cite{BHR12} is a tuple $(\Gb, \En, \Ev, \De)$. On a circuit
$F$, $\Gb(1^{k}, F)$ outputs a garbled circuit $\hat F$, encoding information
$e$ and decoding information $d$. $\En(e, x)$ maps a plaintext input to garbled
input labels, $\Ev(\hat F, X)$ evaluates, and $\De(d, Y)$ decodes. We use the
three standard properties: \emph{correctness}
($\De(d,\Ev(\hat F,\En(e,x))) = F(x)$), \emph{privacy}, and \emph{authenticity}.

Authenticity, verbatim in structure from~\cite{ZRE15}: given $(\hat F, X)$
alone, no adversary can produce $\tilde Y \ne \Ev(\hat F, X)$ with
$\De(d, \tilde Y) \ne \bot$ except with negligible probability. The quantifier is
the point. The adversary receives the garbled circuit and one encoded input, and
nothing else.

\subsection{Half-gates with free-XOR}
\label{sec:zre}

We work with the half-gates scheme of~\cite{ZRE15} with the free-XOR
optimisation of~\cite{KS08}, which is what the BitVM Alliance implementation
uses. Writing $\Hh$ for a circular-correlation-robust hash~\cite{GKWY20}:

\begin{align*}
\Gb:\quad & R \leftarrow\!\!\text{\$}\ \bit^{k-1}1, \qquad
            W_i^{0} \leftarrow\!\!\text{\$}\ \bit^{k}\ \text{ for inputs},\qquad
            W_i^{1} = W_i^{0} \oplus R \\
          & \text{XOR gate: } W_i^{0} = W_a^{0} \oplus W_b^{0}
            \qquad\text{(no ciphertexts)} \\
          & \text{AND gate: } (W_i^{0}, T_{Gi}, T_{Ei}) \leftarrow
            \textsf{GbAnd}(W_a^{0}, W_b^{0}) \\
          & \text{output: } d_i = \lsbf(W_i^{0}) \\[0.3em]
\En:\quad & X_i = e_i \oplus x_i R \\[0.3em]
\Ev:\quad & \text{XOR: } W_i = W_a \oplus W_b \\
          & \text{AND: } s_a = \lsbf(W_a),\ s_b = \lsbf(W_b), \\
          & \qquad W_{Gi} = \Hh(W_a, j) \oplus s_a T_{Gi}, \qquad
            W_{Ei} = \Hh(W_b, j') \oplus s_b (T_{Ei} \oplus W_a), \\
          & \qquad W_i = W_{Gi} \oplus W_{Ei} \\[0.3em]
\De:\quad & y_i = d_i \oplus \lsbf(Y_i)
\end{align*}

Two structural facts drive everything below. First, every wire's two labels
differ by the same global offset $R$, so any gate combining labels by XOR is
$\mathbb{F}_2$-affine in them. Second, the scheme as published provides
correctness and privacy but, in the authors' own statement, does not provide
authenticity; authenticity requires a modified decoding procedure.

\subsection{BitVM3-core}

BitVM3-core~\cite{BitVM3} instantiates the above with $k = 128$. The prover
garbles the SNARK verifier, publishes $\hat F$ and $d$ off-chain, and posts an
\textsf{Assert} transaction whose witness carries one $65$-byte Schnorr adaptor
signature per $8$-bit digit of the proof. By extractability of that signature
scheme, any observer recovers $\En(e,\pi)$ from the chain. Writing $L^{*}$ for
the output label decoding to \textsc{False}, the \textsf{Assert} output carries a
hash lock on $\Hh(L^{*})$; a challenger who obtains $L^{*}$ by evaluation spends
it and takes the deposit.

Note that $\Hh(Z^{\textsc{True}})$ is never published. In the original design it
need not be: the operator wins by the \emph{absence} of a disproof, so no party
ever has to recognise success.

\section{Evaluation-binding}

\begin{definition}[Evaluation-binding]
\label{def:evb}
A garbling scheme $(\Gb,\En,\Ev,\De)$ is \emph{evaluation-binding} for a circuit
$F$ if for every PPT adversary $\mathcal{A}$ given
$(\hat F, e, d) \leftarrow \Gb(1^{k}, F)$, the probability that $\mathcal{A}$
outputs a label vector $L$ with $\De(d, \Ev(\hat F, L)) = \textsc{True}$, when
$F$ is unsatisfiable, is negligible.
\end{definition}

Two remarks on the shape of the definition.

First, it must be stated over \emph{outcomes} and not over the form of $L$. The
tempting statement, that $L$ must equal $\En(e,x)$ for some satisfying $x$, is
false; see Attack~\ref{atk:cancel}.

Second, the property is not new to this paper in substance. Glock~\cite{Glock}
states its security goal as: the evaluator learns the output label \emph{if and
only if} the garbler authenticated an input on which the circuit evaluates to
that output. The ``only if'' direction is evaluation-binding. Every existing
construction discharges it by the on-chain reveal rather than by proving
anything about the garbling scheme. What is new here is the observation that the
obligation exists, and what happens when one tries to meet it directly.

\begin{observation}
Evaluation-binding is not implied by correctness, privacy or authenticity.
Correctness quantifies only over valid encodings. Privacy is irrelevant, and
indeed the schemes used in this setting are deliberately privacy-free. And the
authenticity experiment withholds $e$, $d$ and $R$ from the adversary, whereas
the garbler holds all three.
\end{observation}

\section{Attacks}
\label{sec:attacks}

\begin{attack}[Native decoding admits trivial forgery]
\label{atk:lsb}
Under the decoding of Section~\ref{sec:zre}, $\De(d, Y) = d \oplus \lsbf(Y)$ is
a single XOR against one bit of the label. Every $k$-bit string decodes to a
valid output bit; there is no $\bot$. A malicious garbler therefore samples
arbitrary $L$, evaluates locally, and resamples if the result decodes to
\textsc{False}. The expected number of trials is two.
\end{attack}

Our implementation forges in $1$, $1$ and $3$ trials at $k = 32, 64, 128$
respectively. The repair is standard~\cite{ZRE15}: commit to both output labels,
$d = (\Hh(Z^{0}), \Hh(Z^{1}))$, and return $\bot$ on a mismatch. All subsequent
attacks assume this repair.

\begin{attack}[One-sided decoding commitment]
\label{atk:onesided}
BitVM3-core publishes $\Hh(L^{*})$ and not $\Hh(Z^{\textsc{True}})$. A party
evaluating off-chain can therefore recognise failure and nothing else:
\textsc{True} and garbage are the same observation. Any protocol that asks a
challenger to act on a successful evaluation is unimplementable as specified.
\end{attack}

The repair is $32$ bytes at setup. We flag it because it is invisible while the
on-chain reveal is present, and becomes load-bearing the moment it is removed.

\begin{attack}[Encoding-form binding is false]
\label{atk:cancel}
Let $S$ be any set of input wires whose parity into every application of $\Hh$
is even; two wires meeting at a common XOR gate suffice. For any
$\rho \notin \{0, R\}$ set $L_m = W_m^{0} \oplus \rho$ for $m \in S$ and
$L_m = W_m^{x_m}$ otherwise. The $\rho$ cancels at the XOR, no hash is ever
applied to a perturbed label, and the evaluation is bit-identical to the honest
one. Thus $L$ is not $\En(e,x)$ for any $x$, yet decoding succeeds.
\end{attack}

This does not help a cheating prover on its own. It matters because it fixes the
shape of Definition~\ref{def:evb}, and because the mechanism it exposes, free-XOR
rendering regions of the circuit $\mathbb{F}_2$-affine in the labels, is what the
next attack exploits.

\begin{attack}[Birthday attack: the $2^{k/2}$ ceiling]
\label{atk:birthday}
Suppose the output label can be written $Z = g(L_S) \oplus h(L_T)$ for disjoint
input-wire sets $S, T$, which holds whenever the paths from $S$ and from $T$
merge only through XOR gates after their last application of $\Hh$. Then:
tabulate $g(L_S)$ over $2^{k/2}$ garbage assignments to $L_S$; sweep $2^{k/2}$
garbage assignments to $L_T$ looking up $Z^{\textsc{True}} \oplus h(L_T)$. A
collision is expected and yields $L$ with $\Ev(\hat F, L) = Z^{\textsc{True}}$.
\end{attack}

Modelling $\Hh$ as a random oracle, an unstructured search for the target costs
$2^{k}$. The split reduces it to $2^{k/2}$, and van Oorschot--Wiener parallel
collision search~\cite{vOW99} makes that memory-light and perfectly
parallelisable. For the deployed parameter $k=128$ the ceiling is $2^{64}$.

At $10^{9}$ hash evaluations per core-second, $2^{64}$ is roughly $585$
core-years, or about three weeks on $10{,}000$ cores. Whether that is alarming
depends on the value secured; Clementine~\cite{Clementine} has processed on the
order of $150$~BTC.

\begin{attack}[No partial repair to the soundness proof]
\label{atk:proof}
BitVM3's soundness proof derives $L_\pi = \En(e,\pi)$ from extractability and
only then applies correctness. Removing the on-chain reveal deletes the premise
of the second step rather than weakening its conclusion. There is no
intermediate position in which a weaker on-chain reveal yields a weaker but
still sufficient guarantee.
\end{attack}

\section{Implementation}
\label{sec:impl}

We implement the half-gates scheme of Section~\ref{sec:zre} in Rust,
parameterised by label length $k$, following~\cite{ZRE15} Figure~2 line by line.
We do not use the BitVM Alliance implementation~\cite{GSV} here because it fixes
$k = 128$, placing Attack~\ref{atk:birthday} beyond demonstration. Before any
attack runs, the garbling is validated exhaustively against plaintext evaluation
on every input of every test circuit at $k \in \{16,32,64,128\}$.

The test circuit for Attack~\ref{atk:birthday} is
$\textsf{out} = (a \wedge b) \oplus (c \wedge d)$: two AND gates and four
inputs, the minimal circuit exhibiting the split. We hold $L_b, L_d$ at honest
labels and vary $L_a, L_c$ over garbage.

\begin{table}[t]
\centering
\begin{tabular}{rrrrcc}
\toprule
$k$ & $\Hh$ evaluations & $2^{k/2}$ & $2^{k}$ & no valid encoding & decodes \\
\midrule
24 & $20{,}722$ & $4.10\!\times\!10^{3}$ & $1.68\!\times\!10^{7}$ & yes & \textsc{True} \\
28 & $105{,}924$ & $1.64\!\times\!10^{4}$ & $2.68\!\times\!10^{8}$ & yes & \textsc{True} \\
32 & $280{,}360$ & $6.55\!\times\!10^{4}$ & $4.29\!\times\!10^{9}$ & yes & \textsc{True} \\
36 & $1{,}170{,}344$ & $2.62\!\times\!10^{5}$ & $6.87\!\times\!10^{10}$ & yes & \textsc{True} \\
40 & $4{,}349{,}874$ & $1.05\!\times\!10^{6}$ & $1.10\!\times\!10^{12}$ & yes & \textsc{True} \\
44 & $24{,}420{,}688$ & $4.19\!\times\!10^{6}$ & $1.76\!\times\!10^{13}$ & yes & \textsc{True} \\
\bottomrule
\end{tabular}
\caption{Attack~\ref{atk:birthday}, measured. ``No valid encoding'' is an
exhaustive check that the forged vector equals $\En(e,x)$ for none of the $2^{n}$
inputs. Decoding uses the committed variant, so \textsc{True} means the
evaluation landed exactly on $Z^{\textsc{True}}$.}
\label{tab:birthday}
\end{table}

\begin{table}[t]
\centering
\begin{tabular}{rrrr}
\toprule
$k$ & birthday $\Hh$ evals & naive $\Hh$ evals & speedup \\
\midrule
20 & $5{,}702$ & $1{,}160{,}196$ & $203\times$ \\
24 & $20{,}722$ & $36{,}204{,}000$ & $1{,}747\times$ \\
28 & $105{,}924$ & $345{,}161{,}504$ & $3{,}259\times$ \\
\bottomrule
\end{tabular}
\caption{Head to head against unstructured search for the same target at the
same $k$. The speedup itself grows as $2^{k/2}$.}
\label{tab:headtohead}
\end{table}

Table~\ref{tab:birthday} reports the attack, Table~\ref{tab:headtohead} the
comparison against unstructured search. Fitting the measured work against $k$
over $[24,44]$ gives a scaling exponent of $\mathbf{1.02}$, against $1.00$
predicted by the birthday bound. Every row of Table~\ref{tab:birthday} is a
label vector that is the encoding of no input whatsoever and which an honest
evaluator accepts as a proof of \textsc{True}.

\section{Does the deployed circuit have the structure?}
\label{sec:measure}

Attack~\ref{atk:birthday} needs a split. Section~\ref{sec:impl} demonstrates it
on a circuit built to have one. This section asks whether the circuit BitVM3-core
actually garbles has one.

\subsection{Method}

We implement a custom evaluation mode for the BitVM Alliance streaming circuit
builder~\cite{GSV} and run the real \texttt{groth16\_verify} circuit end to end.
Each wire carries two pieces of information.

\emph{A dependency sketch.} Input wire $i$ is seeded with bit $i \bmod 64$ and
every gate ORs its operands, so a wire's sketch records which of $64$ residue
classes of input wires it depends on.

\emph{A recorded tail.} The last $2^{21}$ gates are kept with their wire
identifiers and sketches, which lets the \emph{linear frontier} of the output be
reconstructed exactly offline: from the output wire, step back through free
gates only, stopping at non-free gates or circuit inputs, with even-parity atoms
cancelling as the label algebra requires. The output label is the XOR of its
frontier atoms' labels, so a split at the output exists only if the frontier has
at least two atoms with separable dependencies.

\subsection{Results}

\begin{table}[t]
\centering
\begin{tabular}{lr}
\toprule
input wires & $3{,}302$ \\
gates executed & $10{,}338{,}696{,}819$ \\
non-free gates & $2{,}714{,}835{,}821$ \\
output sketch popcount & $64/64$ \\
\midrule
backward steps through free gates & $1$ \\
frontier atoms (odd parity) & $\mathbf{1}$ \\
unresolved wires & $0$ \\
\midrule
top gate operand A / B popcount & $64/64$, $64/64$ \\
mutually exclusive residue classes & none \\
\midrule
non-free gates in recorded tail & $638{,}658$ \\
\quad with separably controllable operands & $\mathbf{0}$ \\
\bottomrule
\end{tabular}
\caption{One full pass over the BN254 Groth16 verifier circuit.}
\label{tab:measure}
\end{table}

Table~\ref{tab:measure} gives the result. Three findings of differing strength.

\begin{enumerate}[leftmargin=*,itemsep=2pt]
  \item \textbf{Exact.} The output's linear frontier is a single atom. The
    backward walk terminated after one step with nothing unresolved, so the
    output wire is produced directly by a non-free gate with no XOR above it.
    The top-level decomposition Attack~\ref{atk:birthday} requires does not
    exist. This is not statistical.
  \item \textbf{Evidence.} The top non-free gate's operands each depend on
    inputs from all $64$ residue classes with no mutually exclusive class, so
    they are not separably controllable at this resolution.
  \item \textbf{Evidence.} Across the $638{,}658$ non-free gates in the last
    $2^{21}$ gates of the circuit, none has separably controllable operands.
\end{enumerate}

\begin{claim}
Attack~\ref{atk:birthday} does not apply to the deployed BN254 Groth16 verifier
circuit.
\end{claim}

The reason is structural: the Groth16 verifier terminates in an equality check
over the whole final-exponentiation result, so every wire near the output
depends on essentially every input and there is nothing to vary independently.

\subsection{Scope}

The recorded tail is $2^{21}$ of $1.03 \times 10^{10}$ gates, the last
$0.02\%$. Finding~1 is exact and unconditional. Findings~2 and~3 hold at
$64$-class resolution, and two cones could in principle be disjoint while both
touching all $64$ classes; an exact analysis would carry a $3302$-bit set per
live wire, roughly $66$~MB resident, and is future work. Separately, our
non-free gate count of $2{,}714{,}835{,}821$ differs from the upstream gate
counter's $2{,}715{,}041{,}234$ by $205{,}413$, which we believe to be gates
whose output wire is unreachable and which the upstream counter includes; this
is unverified.

\section{Discussion}

\subsection{From break to design requirement}

The deployed circuit satisfies the structural condition, and no one chose that.
It follows from Groth16's output convention, not from a decision. It is not
obviously preserved under a change of proof system, a change of circuit
compiler, or a change of output encoding. A hash-based verifier terminating in
an XOR-tree accumulator, which is where the post-quantum direction for these
protocols points, is a natural place for it to fail.

We therefore propose that evaluation-binding, and the structural condition that
supports it, be stated and checked rather than inherited. The tool accompanying
this paper decides the condition in one pass over the circuit.

\subsection{Is $k = 128$ the right parameter here?}

The $128$-bit label length is inherited from two-party computation, where the
garbler cannot choose the evaluator's input. BitVM-style protocols are not that
setting whenever labels reach an evaluator without on-chain authentication. If
Attack~\ref{atk:birthday} applies to some future circuit, restoring $128$-bit
security requires $k = 256$, which doubles the garbled circuit. For the
verifier measured here that is $40.5$~GB becoming $81$~GB, in a design where
off-chain size is the binding constraint. The parameter deserves explicit
justification in this setting rather than inheritance.

\subsection{Relation to prior attacks}

Futoransky, Larotonda and Barbara~\cite{FairgateNote} broke an earlier
RSA-based garbling proposal for this same application with a circuit containing
two AND gates. Their attack is evaluator-side and exploits multiplicative
structure in $\mathbb{Z}_N$, which half-gates with a correlation-robust hash
does not have, so it does not transfer. The shape is instructive nonetheless:
in both cases the claim is that a party cannot produce labels it should not be
able to produce, the intuition was that structure made it hard, and the
structure was exploitable.

\section{Open problems}

\begin{enumerate}[leftmargin=*,itemsep=2pt]
  \item Prove or refute evaluation-binding for half-gates under circular
    correlation robustness~\cite{GKWY20}, for circuits satisfying the structural
    condition. Attack~\ref{atk:birthday} bounds what any such proof can achieve
    at $2^{k/2}$ in general.
  \item Give a decision procedure for the structural condition that is exact
    rather than sketch-based, and cheap enough to run in continuous integration.
  \item Characterise which circuit output conventions guarantee the condition by
    construction, so that it can be enforced by the compiler rather than audited
    after the fact.
  \item Determine whether arithmetic garbling schemes
    (\cite{Glock,DutyFreeBits}) admit an analogous garbler-side binding gap, and
    whether their algebraic structure makes it easier or harder to exploit.
\end{enumerate}

\section*{Artifacts}

All code and measurements are available at
\url{https://github.com/aryaethn/garbler-binding}. The exploit is a single Rust
file with one dependency and runs in about ten seconds. The circuit analysis
requires one pass over the BN254 Groth16 verifier, roughly ten minutes on a
laptop core.

\section*{Acknowledgements}

This work was carried out with assistance from Claude (Anthropic)
for literature reconciliation and implementation.

\bibliographystyle{plain}
\bibliography{refs}

\end{document}
